\section{Related Work}
\label{sec:related}

\authornote{\textbf{Diffusion posterior sampling for MRI.} Score-based CS-MRI
\citep{jalal2021robust,chung2022scoremri}; priors trained on corrupted data
\citep{aali2025ambient}; frequency-split diffusion processes that \emph{require
training}, namely HFS-SDE \citep{cao2024hfssde} and CDiffMR \citep{huang2023cdiffmr}.
Benchmarking: InverseBench includes multi-coil fastMRI with real raw measurements and
tunes each sampler per task \citep{zheng2025inversebench}; we follow its tuning
protocol. One sentence on where we differ: we are training-free and change only the
DC step.}

\authornote{\textbf{Beyond isotropic posterior covariance.} \citet{peng2024optimal} show
that DPS and \pigdm{} share an isotropic Gaussian approximation and replace it with a
covariance that is diagonal in a wavelet basis and learned or estimated offline. TMPD
\citep{boys2024tmpd} and MMPS \citep{rozet2024mmps} use the denoiser Jacobian (MMPS is
also evaluated on fastMRI knee). CA-DPS \citep{hamidi2024cadps} estimates a diagonal
covariance online. \citet{schioppa2026covariance} use a Fourier-structured covariance
for unconditional sampling. \emph{Difference (state it precisely):} our covariance is
diagonal in the basis that diagonalises the MRI operator, is calibrated per scan from
the ACS region with no extra network evaluations, and is evaluated on raw multi-coil
and low-field data.}

\authornote{\textbf{Frequency scheduling and adaptive DC.} \citet{tian2026frequency}
expose measurement frequencies according to the diffusion noise level and return
progressively to full-band consistency, on natural images. Our frequency gate is the
k-space instance of this idea, and we compare against their schedule
(\tableref{tab:ablation}). SMRD \citep{ozturkler2023smrd} adapts a \emph{scalar} DC
weight with SURE and stops early. Measured-subspace consistency \citep{lee2026msc} is
a plug-in range/null-space correction for MRI samplers. It is complementary to ours,
because it does not reweight frequencies. Blind or noise-adaptive samplers estimate a
scalar noise level or a blur kernel
\citep{laroche2024fastdiffem,xia2026nanhmc}.}

\authornote{\textbf{Noise in MRI.} Thermal noise is white per channel and correlated
across coils; prewhitening with the coil noise covariance is standard
\citep{roemer1990phased,pruessmann1999sense,kellman2005snr}. Magnitude images follow
Rician or noncentral-$\chi$ statistics \citep{gudbjartsson1995rician,constantinides1997snr}.
Keep this to 2--3 sentences. It justifies the white-noise likelihood in
\equationref{eq:lik} and the magnitude-domain evaluation caveat.}
